\documentclass[a4paper,10pt]{article}
\usepackage[dutch]{babel}
\usepackage{amsmath}
\usepackage{amssymb}
\usepackage{bbm}
\usepackage[amssymb]{SIunits}%voor grootheden
\usepackage[pdftex]{graphicx}%voor figuren
\usepackage[margin=1.9cm,nohead]{geometry}
\usepackage{eurosym} %voor het euroteken
\usepackage[version=3]{mhchem} %voor chemische reacties


%hoofding voor de taak, niks aan veranderen
\newcommand{\hoofding}[6]{
\begin{flushleft}
\includegraphics[height=2.5cm]{logo_hhvmb.jpg}
\end{flushleft}
\vspace{-2.8cm}
\hspace{4cm}
\parbox{8cm}{Naam : #1\\Klas : #2\\
             Nummer : #3\\Datum : #4\\
             Vak : #5\\Leerkracht : #6\\}
\hspace{2.3cm}   
\begin{tabular}{|p{1.7cm}|}
							\hline
							\vspace{1cm}\\
							\hline
\end{tabular}

{\parindent=0pt \hrulefill} %lege regel invoegen want moet als paragraaf gezien worden

\vspace{5mm}}
%einde hoofding van toetsen


% figuur invoegen
\newcommand{\figuur}[2]{\includegraphics[width=#1]{#2}} 



%hiermee is er geen insprong bij een nieuwe alinea
\setlength{\parindent}{0pt}
%de witruimte tussen twee alinea's wordt iets groter
\setlength{\parskip}{1ex plus 0.5ex minus 0.2ex}









\begin{document}

%naam,klas,nr,datum,vak,leerkracht
\hoofding{*}{*}{*}{*}{vak}{*}

					
%hieronder de titel typen
\begin{center}
\LARGE{\textbf{Titel}}
\end{center}


%paragraaftitels
\section{Paragraaf}

\subsection{paragraaf}

Begin van de verslagtekst.




%informatie om een figuur in te voegen
Met \verb:\figuur{8cm}{tabel}: voeg je de onderstaande figuur in.
\begin{center}
\figuur{8cm}{tabel}
\end{center}
	
 Met \verb:$$\rho=\frac{m}{V}$$:  voeg je
 de onderstaande vergelijking in.  
$$\rho=\frac{m}{V}$$
  
 Met \verb:$$v=\unit{2,3\cdot 10^4}{\meter\per\second}$$:  voeg je
 de onderstaande grootheid in. 
$$v=\unit{2,3\cdot 10^4}{\meter\per\second}$$
  
  
 Met \verb:\ce{HgI2 ->[I-] [Hg^{II}I4]^2-}:  voeg je
 de onderstaande reactie in.
 
 \ce{HgI2 ->[I-] [Hg^{II}I4]^2-}  
   
    
 \end{document} 


\end{document}